% Ad Familiares 8.13 (ad-familiares-08-13)
% Source: https://raw.githubusercontent.com/PerseusDL/canonical-latinLit/master/data/phi0474/phi056/phi0474.phi056.perseus-lat1.xml
% Fetched by scripts/fetch_latin.py

% Dateline (Perseus): Scr. Romae ex. m. Mai. aut in. Iun. a. 704 (50) .

\ciceroLetterOpener{CAELIVS CICERONI S.}

\ciceroSection{1} gratulor tibi adfinitate viri medius fidius optimi ; nam hoc ego de illo existimo. cetera porro, quibus adhuc ille sibi parum utilis fuit, et aetate iam sunt decussa, et consuetudine atque auctoritate tua, pudore Tulliae, si qua restabunt, confido celeriter sublatum in ; non est enim pugnax in vitiis neque hebes ad id quod melius sit intellegendum. deinde , quod maximum est, ego illum valde amo.

\ciceroSection{2} voles , Cicero, Curionem nostrum lautum intercessionis de provinciis exitum habuisse ; nam cum de intercessione referretur, quae relatio fiebat ex senatus consulto, primaque M. Marcelli sententia pronuntiata esset, qui agendum cum tribunis pl. censebat, frequens senatus in alia omnia iit. stomacho est scilicet Pompeius Magnus nunc ita languenti, ut vix id quod sibi placeat reperiat. transierant illuc, rationem eius habendam qui neque exercitum neque provincias traderet. quem ad modum hoc Pompeius laturus sit, cum cognoscam ; quidnam rei publicae futurum sit, si †aut non curet, vos senes divites videritis. Q. Hortensius, cum has litteras scripsi, animam agebat.
